\subsection{子环；理想；商环；环的积}

若 H 是拓扑环 A 的子环，则由 A 的拓扑在 H 上诱导的拓扑与 H 的环结构相容。这样在 H 上定义的拓扑环结构称为由 A 的拓扑环结构所诱导。

命题 5. 设 H 是拓扑环 A 的稠密子环，并设 K 是 H 的子环（相应地，左理想、右理想、双边理想）。则 K 在 A 中的闭包 K 是 A 的子环（相应地，左理想、右理想、双边理想）。

证明与 § 2 第 3 节命题 8 的证明相同：例如若 K 是 H 中的左理想，则映射 \((z, x) \to zx\) 在 \(A \times A\) 上连续，且把 \(H \times K\) 映入 K；因而它把 \(A \times K = \overline{H} \times K\) 映入 \(\overline{H}\)。

设 H 是拓扑环 A 中的双边理想。用与商群情形同样的论证，我们看到 A 的拓扑由关系 \(x - y \in H\) 作出的商与 A/H 的环结构相容。特别地，若 A 不是豪斯多夫的，则 \(\{o\}\) 在 A 中的闭包 N 由命题 5 是一个闭双边理想；商环是豪斯多夫的（§ 2，第 5 节，命题 13），称为与 A 相伴的 Hausdorff 环。

设 \((\mathbf{A}_t)_{t\in \mathbf{I}}\) 是拓扑环的一个族。在集合 \(\mathbf{A} = \prod_{t\in \mathbf{I}}\mathbf{A}_t\) 上，诸 \(\mathbf{A}_t\) 的拓扑之积与诸 \(\mathbf{A}_t\) 的环结构之积相容（与积群情形的证明相同）；这样定义的拓扑环 \(\mathbf{A}\) 称为诸拓扑环 \(\mathbf{A}_t\) 的积。

